activity (measured in IU/L) after 90 min of exercise in 15 healthy young men (Totsuka, M., et al., Break point of serum creatine kinase release after endurance exercise. \url{http://jap.physiology.org/cgi/content/full/93/4/1280}):

\begin{center}
\begin{tabular}{cccccccccc}
1112 & 722 & 689 & 251 & 196 & 185 & 128 & 102 & 166 & 178 \\
775 & 694 & 514 & 244 & 208 & & & & &
\end{tabular}
\end{center}

Construct a 95\% CI for the mean peak CK activity.

\begin{description}
\item[5.5.36.] A random sample of 20 observations gave the following summary statistics: $\sum x_i = 234$ and $\sum x_i^2 = 3048$. Assuming that the data may be looked upon as a random sample from a normal population, construct a 95\% CI for the actual average, $\mu$.
\item[5.5.37.] Let a random sample of size 17 from a normal population for which both mean $\mu$ and variance $\sigma^2$ are unknown yield $\bar{x} = 3.12$ and $s^2 = 1.04$. Determine a 99\% CI for $\mu$.
\item[5.5.38.] A random sample from a normal population yields the following 25 values:
\end{description}

\begin{center}
\begin{tabular}{cccccccccc}
90 & 87 & 121 & 96 & 106 & 107 & 89 & 107 & 83 & 92 \\
117 & 93 & 98 & 120 & 97 & 109 & 78 & 87 & 99 & 79 \\
104 & 85 & 91 & 107 & 89 & & & & &
\end{tabular}
\end{center}

\begin{enumerate}
\item[(a)] Calculate an unbiased estimate $\bar{\theta}$ of the population mean.
\item[(b)] Give an approximate 99\% CI for the population mean.
\end{enumerate}

\begin{description}
\item[5.5.39.] The following are random data from a normal population.
\end{description}

\begin{center}
\begin{tabular}{cccccccccc}
3.3 & 3.3 & 4.7 & 2.6 & 6.4 & 4.7 & 1.7 & 4.5 & 5.0 & 3.0
\end{tabular}
\end{center}

Construct a 98\% CI for the population mean $\mu$.

\begin{description}
\item[5.5.40.] The following data represent the rates (micrometers per hour) at which a razor cut made in the skin of anesthetized newts is closed by new cells:
\end{description}

\begin{center}
\begin{tabular}{cccccccccc}
28 & 20 & 21 & 39 & 32 & 23 & 18 & 31 & 14 & 23 \\
18 & 22 & 28 & 24 & 33 & 12 & 23 & 21 & 25 & 25
\end{tabular}
\end{center}

\begin{enumerate}
\item[(a)] Can we say that the data are approximately normally distributed?
\item[(b)] Find a 95\% CI for population mean $\mu$ for the new cells to close a razor cut made in the skin of anesthetized newts.
\item[(c)] Find a 99\% CI for $\mu$.
\item[(d)] Is the 95\% CI wider or narrower than the 99\% CI? Briefly explain why.
\end{enumerate}

\begin{description}
\item[5.5.41.] For a particular car, when the brake is applied at 62 mph, the following data give stopping distance (in feet) for 10 random trials on a dry surface (source: \url{http://www.nhtsa.dot.gov/cars/testing/brakes/b/pdf}):
\end{description}

\begin{center}
\begin{tabular}{cccccc}
146.9 & 148.4 & 149.4 & 148.6 & 150.3 \\
147.5 & 147.5 & 149.3 & 148.4 & 145.5
\end{tabular}
\end{center}

\begin{enumerate}
\item[(a)] Can we say that the data are approximately normally distributed?
\item[(b)] Find a 95\% CI for population mean stopping distance, $\mu$.
\end{enumerate}

\begin{description}
\item[5.5.42.] A pharmaceutical company tested a new medicine to be marketed for the treatment of a particular type of virus. To obtain an estimate of the mean recovery time, this medicine was tested on 15 volunteer patients, and the recovery time (in days) was recorded. The following data were obtained:
\end{description}

\begin{center}
\begin{tabular}{cccccccccc}
8 & 17 & 10 & 6 & 34 & 11 & 13 & 6 & 9 & 8 \\
19 & 4 & 12 & 17 & 7 & & & & &
\end{tabular}
\end{center}

\begin{enumerate}
\item[(a)] Obtain a 95\% CI estimate of the mean recovery.
\item[(b)] What assumptions do we need to make? Test for these assumptions.
\end{enumerate}
